\documentclass[12pt]{article}
\usepackage[utf8]{inputenc}
\usepackage[T1]{fontenc}
\usepackage{hyperref}
\title{Waddar Language Resources for K-12 Curriculum Development}
\author{Compiled by Claude Sonnet 5 (claude-sonnet-5) for Andr\'as Kornai}
\date{}
\begin{document}
\maketitle

\section{Introduction}
Waddar (Vadari, Odde; ISO 639-3 wbq; Glottocode wadd1237) is a Dravidian language of India. It is usually placed near Telugu, although that placement is questioned. The 2011 census figure found in search results is about 200,000 speakers. Waddar speakers also use Telugu and Kannada. The search found very little material specific to Waddar. No school curriculum, primer, textbook, grammar or dictionary was found.

\section{Primary Educational Resources}
No Waddar-language primers, textbooks or school programmes were found in search results.

\section{Linguistic and Reference Materials}
\begin{itemize}
\item Waddar language, Wikipedia overview: \url{https://en.wikipedia.org/wiki/Waddar_language}
\item ASJP Waddar wordlist, based on Grierson's 1906 Linguistic Survey of India: \url{https://asjp.clld.org/languages/WADDAR}
\item Max Planck Institute numerals page for Waddara (Vadari): \url{https://lingweb.eva.mpg.de/channumerals/Waddar.htm}
\item Burrow and Emeneau, A Dravidian Etymological Dictionary (a general comparative work; Waddar entries were not confirmed): \url{https://dsal.uchicago.edu/dictionaries/burrow/new_frontmatter.html}
\end{itemize}

\section{Religious and Community Resources}
Joshua Project has a profile of the Vaddar (Hindu traditions) people group in India: \url{https://www.joshuaproject.net/index.php/people_groups/18288/IN}. No Waddar scripture translations or community organisations were found.

\section{Digital and Audio Resources}
No Waddar audio, video or corpus resources were found. The FontSpace language page for wbq (Telugu script) was listed, which suggests Telugu script is the likely writing system: \url{https://www.fontspace.com/languages/info/wbq-telu}

\section{Curriculum Development Context}
The language is not known to be taught in schools, and the evidence found is thin. A curriculum would need new primary documentation. A reasonable first step is the 1906 Linguistic Survey of India material and fieldwork with the community. Telugu and Kannada are the regional school languages.
\end{document}